\documentclass[12pt]{ximera}
\title{Homework 5: Concepts of Probability and Odds}
\begin{document}
\begin{abstract}
Section 7.4: probabilities of ``or'' events for a pair of dice, inclusion-exclusion
with probabilities, and converting between probability and odds --- including why a
roulette wheel's house odds differ from its true odds.
\end{abstract}
\maketitle

\section*{Section 7.4: Concepts of Probability}

\begin{problem}
Suppose a d4 and a d6 are rolled (that is, a $4$-sided die and a $6$-sided die).
Determine the probabilities of the following events. Enter fractions in lowest terms.

\begin{prompt}
\textbf{(a)} The d4 shows a $3$ or the sum is $7$.

$\answer{3/8}$

\begin{hint}
There are $4\cdot 6=24$ equally likely outcomes. Count each event, then subtract the
outcomes you counted twice.
\end{hint}

\begin{feedback}
Of the $24$ equally likely outcomes:
\begin{itemize}
\item the d4 shows $3$ in $6$ of them (one for each face of the d6);
\item the sum is $7$ in $4$ of them: $(1,6)$, $(2,5)$, $(3,4)$, $(4,3)$;
\item both happen in $1$: $(3,4)$.
\end{itemize}
By inclusion-exclusion, $6+4-1=9$ outcomes qualify, so the probability is
$\frac{9}{24}=\frac38$.
\end{feedback}
\end{prompt}

\begin{prompt}
\textbf{(b)} The d6 shows a $5$ or the sum is $10$.

$\answer{5/24}$

\begin{feedback}
The d6 shows $5$ in $4$ outcomes (one for each face of the d4). A sum of $10$ needs the
largest faces, and happens only as $(4,6)$ --- $1$ outcome. The two events cannot
happen together, since a d6 showing $5$ would need a d4 showing $5$. So $4+1=5$
outcomes qualify, and the probability is $\frac{5}{24}$.
\end{feedback}
\end{prompt}
\end{problem}

\begin{problem}
Use inclusion-exclusion or a Venn diagram to solve the following.

If $P(E)=0.26$, $P(F)=0.41$, and $P(E\cap F)=0.14$, find each probability.

\begin{prompt}
\textbf{(a)} $P(E\cup F) = \answer[tolerance=0.001]{0.53}$

\begin{feedback}
\[ P(E\cup F)=P(E)+P(F)-P(E\cap F)=0.26+0.41-0.14=0.53. \]
\end{feedback}
\end{prompt}

\begin{prompt}
\textbf{(b)} $P\!\left(\olsi{E}\cap F\right) = \answer[tolerance=0.001]{0.27}$

\begin{hint}
$\olsi{E}\cap F$ is the part of $F$ lying outside $E$. In a Venn diagram, that is the
$F$ circle with the overlap removed.
\end{hint}

\begin{feedback}
The event $F$ splits into the part inside $E$ and the part outside it, so
\[ P\!\left(\olsi{E}\cap F\right)=P(F)-P(E\cap F)=0.41-0.14=0.27. \]
\end{feedback}
\end{prompt}
\end{problem}

\begin{problem}
Consider the following questions about odds.

\begin{prompt}
\textbf{(a)} According to the U.S.\ Bureau of Transportation Statistics, in 2020 about
$85\%$ of the flights on U.S.\ carriers were on time (within $14$ minutes of the
scheduled time). By that estimate, what were the odds of a flight being on time in
2020? Reduce your answer.

The odds are $\answer{17}$ to $\answer{3}$.

\begin{hint}
Odds compare the chance the event happens with the chance it does not.
\end{hint}

\begin{feedback}
On time has probability $85\%$ and not on time $15\%$, so the odds are
\[ 85:15 = 17:3. \]
\end{feedback}
\end{prompt}

In a common form of roulette, a ball is spun around a wheel with $38$ slots of equal
size and eventually lands in one of the slots. The slots are numbered $0$, $00$, and $1$
through $36$.

One wager that can be made in roulette is to bet that the ball will land on a
particular number between $1$ and $36$. The house odds for betting on a single number
are $35$ to $1$ --- meaning that if someone loses the wager then their stake is lost, and
if they win then they earn a profit of $35$ times their stake.

House odds are set to differ from the true odds, so that over time a casino makes a
profit from anyone gambling there.

\begin{prompt}
\textbf{(b)} What are the true odds of losing this type of wager?

The odds are $\answer{37}$ to $\answer{1}$.

\begin{feedback}
Exactly $1$ slot wins and the other $37$ lose, so the true odds of losing are
$37$ to $1$.

The house pays as though the odds were $35$ to $1$. That gap is the casino's edge: a
fair payout for a $1$-in-$38$ chance would be $37$ times your stake, not $35$.
\end{feedback}
\end{prompt}

\begin{prompt}
\textbf{(c)} What is the probability of winning this type of wager? Enter a fraction.

$\answer{1/38}$

\begin{feedback}
One of the $38$ equally likely slots wins, so $P(\text{win})=\frac{1}{38}\approx 0.026$.
Equivalently, odds of $37$ to $1$ against correspond to a probability of
$\frac{1}{37+1}$ of winning.
\end{feedback}
\end{prompt}
\end{problem}

\end{document}
